\clearpage
\VerdierPhysicalPageAnchor{0142}%
\begingroup
\small
\setlength{\parskip}{0pt}%
\setlength{\emergencystretch}{1.5em}%

\VerdierSemanticAnchor{ch2-2-2-11}
\noindent\textbf{推论 2.2.11.}\quad
\begingroup\itshape
\VerdierSemanticAnchor{ch2-2-2-11-setup}%
设 $\mathcal D$ 为三角范畴，
$\mathcal B\subset\mathcal D$ 为满三角子范畴。

\VerdierSemanticAnchor{ch2-2-2-11-a}
\par\noindent a) 商函子
$Q\colon\mathcal D\longrightarrow\mathcal D/\mathcal B$ 的核，是包含
$\mathcal B$ 的最小严格满且饱和的三角子范畴。

\VerdierSemanticAnchor{ch2-2-2-11-b-order-isomorphism}
\par\noindent b) 映射 $\mathcal B\longmapsto S(\mathcal B)$ 把每个
严格满且饱和的三角子范畴对应到它的饱和乘法系 (2.1.8)；它给出一个
有序集同构，把按包含关系排序的严格满且饱和三角子范畴所成的集合，
映到按包含关系排序的 $\mathcal D$ 的饱和乘法系所成的集合上。

\VerdierSemanticAnchor{ch2-2-2-11-b-inverse}
其逆映射把每个饱和乘法系对应到相应局部化函子的核。

\VerdierSemanticAnchor{ch2-2-2-11-c-factorization}
\par\noindent c) 核包含子范畴 $\mathcal B$ 的正合（分别地，上同调）
函子 $F\colon\mathcal D\longrightarrow\mathcal D'$，唯一地通过
$Q\colon\mathcal D\longrightarrow\mathcal D/\mathcal B$ 分解为正合
（分别地，上同调）函子
$G\colon\mathcal D/\mathcal B\longrightarrow\mathcal D'$。

\VerdierSemanticAnchor{ch2-2-2-11-c-morphism-lifting}
核包含子范畴 $\mathcal B$ 的正合（分别地，上同调）函子的态射，
唯一地来自它们的分解之间的态射。
\endgroup

\VerdierSemanticAnchor{ch2-2-2-11-easy-consequences}
这些都是定理 (2.2.6) 的容易推论。

\VerdierSemanticAnchor{ch2-2-2-12}
\smallskip
\noindent\textbf{2.2.12. 关于性质 TRIV 的评论。}
\VerdierSemanticAnchor{ch2-2-2-12-pretriangulated-category}%
把配备一族满足性质 TRI、TRII 与 TRIII 的三角的严格
$\mathbb Z$-范畴称为\emph{前三角范畴}。

\VerdierSemanticAnchor{ch2-2-2-12-induced-notions}
正如对三角范畴那样定义\emph{正合函子}、\emph{上同调函子}与
\emph{乘法系}的概念。

\VerdierSemanticAnchor{ch2-2-2-12-localization}
设 $S$ 为前三角范畴 $\mathcal D$ 中的乘法系 (2.1.4)。局部化范畴
$\mathcal D(S^{-1})$ 具有唯一的前三角范畴结构，使局部化函子为正合
函子。

\VerdierSemanticAnchor{ch2-2-2-12-kernel-subcategory}
$\mathcal D$ 中被局部化函子消去的对象所成的满子范畴，是满前三角
子范畴。

\VerdierSemanticAnchor{ch2-2-2-12-triv-necessity}
实际应用中，还需要知道如何反过来把合适的饱和乘法系关联到一个
前三角子范畴。编者只知道如何利用性质 TRIV 来做到这一点。

\VerdierSemanticAnchor{ch2-2-2-12-triv-role}
因此，这个本质上相当技术性的性质，到目前为止主要用于从定理
(2.2.6) 得到推论 (2.2.11)。

\VerdierSemanticAnchor{ch2-2-3}
\medskip
\noindent\textbf{2.3. 局部化函子的性质。}

\VerdierSemanticAnchor{ch2-2-3-1}
\smallskip
\noindent\textbf{命题 2.3.1.}\quad
\itshape 设 $\mathcal D$ 为三角范畴，$\mathcal B$ 为满三角子范畴。\upshape

\VerdierSemanticAnchor{ch2-2-3-1-p143-boundary}
\endgroup
